%% 
%% Copyright 2007-2020 Elsevier Ltd
%% 
%% This is a modified file from the 'Elsarticle Bundle', which is available under the conditions of the LaTeX Project Public
%% License, either version 1.2 of this license or any later version.
%% ---------------------------------------------

\documentclass[12pt]{elsarticle}

%% For including figures, graphicx.sty has been loaded in
%% elsarticle.cls. If you prefer to use the old commands
%% please give \usepackage{epsfig}

%% The amssymb package provides various useful mathematical symbols
\usepackage{amssymb}
%% The amsthm package provides extended theorem environments
\usepackage{amsthm}

%margin settings
\usepackage[left=2.5cm, right=2.5cm, top=4cm, bottom=3cm, footskip=0.5cm]{geometry}

% suppressing page numbers
\pagenumbering{gobble}

\journal{ }

\begin{document}

\begin{frontmatter}

\title{Title of Your Manuscript}
\date{ }

\author[1,2,3]{Author One}
\author[2,*]{Author Two}
\author[1]{Author Three}

\address[1]{Department of Computing, Research Group, Av. Brasil, 200, São Paulo, SP, Brazil}
\address[2]{School of Science, University of Technology, 2000 J St. NW, Washington DC, 20036, USA}
\address[3]{Center Of Intelligence, Faculty of Sciences, 50 West 4th Street, New York NY, 37895, USA}
\address[*]{Corresponding author: email2@myemail.com}


\begin{abstract}
%% Text of abstract
Lorem ipsum dolor sit amet, consectetur adipiscing elit, sed do eiusmod tempor incididunt ut labore et dolore magna aliqua. Ut enim ad minim veniam, quis nostrud exercitation ullamco laboris nisi ut aliquip ex ea commodo consequat. Duis aute irure dolor in reprehenderit in voluptate velit esse cillum dolore eu fugiat nulla pariatur. Excepteur sint occaecat cupidatat non proident, sunt in culpa qui officia deserunt mollit anim id est laborum.
\end{abstract}

\begin{keyword}
%% keywords here, in the form: keyword \sep keyword
keyword one \sep keyword two \sep keyword three
\end{keyword}

\end{frontmatter}

%% main text
\section{Sample Section Title}
\label{sec:sample1}

%% For citations use: 
%%       \citet{<label>} ==> Jones et al. [21]
%%       \citep{<label>} ==> [21]
%%

Lorem ipsum dolor sit amet, consectetur adipiscing \citep{Fabioetal2013} elit, sed do eiusmod tempor incididunt ut labore et dolore magna \citet{Blondeletal2008} aliqua. Ut enim ad minim veniam, quis nostrud exercitation ullamco laboris nisi ut aliquip ex ea commodo consequat. Duis aute irure dolor in reprehenderit in voluptate velit esse cillum dolore eu fugiat nulla pariatur. Excepteur sint occaecat cupidatat non proident, sunt in culpa qui officia deserunt mollit \citet{Blondeletal2008,FabricioLiang2013} anim id est laborum.

Lorem ipsum dolor sit amet, consectetur adipiscing elit, sed do eiusmod tempor incididunt ut labore et dolore magna aliqua. Ut enim ad minim veniam, quis nostrud exercitation ullamco laboris nisi ut aliquip ex ea commodo consequat. Duis aute irure dolor in reprehenderit in voluptate velit esse cillum dolore eu fugiat nulla pariatur. Excepteur sint occaecat cupidatat non proident, sunt in culpa qui officia deserunt mollit anim id est laborum see appendix~\ref{sec:another}.

\section{Another Section}
\label{sec:another}
Lorem ipsum dolor sit amet, consectetur adipiscing elit, sed do eiusmod tempor section \ref{sec:sample1} incididunt ut labore et dolore magna aliqua. Ut enim ad minim veniam, quis nostrud exercitation ullamco laboris nisi ut aliquip ex ea commodo consequat. Duis aute irure dolor in reprehenderit in voluptate velit esse cillum dolore eu fugiat nulla pariatur. Excepteur sint occaecat cupidatat non proident, sunt in culpa qui officia deserunt mollit anim id est laborum. See Figure~\ref{fig:net}.

\begin{figure}[h]
    \centering
    \includegraphics[width=0.25\textwidth]{net}
    \caption{A network}
    \label{fig:net}
\end{figure}

%% If you have bibdatabase file and want bibtex to generate the
%% bibitems, please use
%%
\bibliographystyle{elsarticle-num-names} 
\bibliography{cas-refs}

%% else use the following coding to input the bibitems directly in the
%% TeX file.

% \begin{thebibliography}{00}

% %% \bibitem[Author(year)]{label}
% %% Text of bibliographic item

% \bibitem[ ()]{}

% \end{thebibliography}
\end{document}

\endinput
%%
%% End of file main.tex'.
